% Generated draft: requires mathematical and source review.
\documentclass[11pt]{article}
\usepackage[margin=25mm]{geometry}
\usepackage{fontspec}
\setmainfont{Latin Modern Roman}
\usepackage{amsmath,amssymb,mathtools,unicode-math}
\setmathfont{Latin Modern Math}
\usepackage{enumitem,needspace,xurl,hyperref}
\setlength{\parindent}{0pt}
\setlength{\parskip}{5pt}
\setlength{\emergencystretch}{3em}
\newcommand{\N}{\mathbb{N}}
\newcommand{\Z}{\mathbb{Z}}
\newcommand{\Q}{\mathbb{Q}}
\newcommand{\R}{\mathbb{R}}
\newcommand{\C}{\mathbb{C}}
\newcommand{\F}{\mathbb{F}}
\newcommand{\A}{\mathbb{A}}
\newcommand{\PP}{\mathbb{P}}
\newcommand{\E}{\mathbb{E}}
\newcommand{\eps}{\varepsilon}
\newcommand{\dd}{\,\mathrm{d}}
\begin{document}
\section*{UnsolvedMath ID: 30002570}
\textbf{Source problem code:} catalogue entry
\subsection{Definitions and mathematical statement}
The source asks for the identification of subcellular processes essential for cancer-cell invasion and for a suitable level of mechanistic detail in a model. Here a mechanistic model means a mathematical description representing biological processes involved in invasion; the source specifies neither a class of such models nor a mathematical definition of invasion. The terms ``essential'', ``realistic'', and ``sufficiently simple'' are criteria to be specified, rather than defined mathematical predicates in the supplied record.

The supplied research question is:
\begin{quote}
Which subcellular processes are essential for cancer-cell invasive behavior, and what level of mechanistic detail is needed for a model that is realistic yet sufficiently simple?
\end{quote}

A precise mathematical formulation remains unresolved. The source does not specify the cancer-cell population or biological setting, the candidate subcellular processes, the invasion observables, a model class, validation data, an adequacy criterion, or a measure of mechanistic complexity. It also does not state whether ``essential'' means necessary for invasion itself or necessary for a model to reproduce specified observations, or whether the requested conclusions should hold across all cancer types and environments or only in a particular setting. No equations, domains, initial or boundary conditions, numerical thresholds, or mathematical quantifiers are supplied. Consequently, the question cannot be converted into a uniquely specified theorem, existence problem, or minimization problem without additional assumptions.

The supplied literature assessment describes models incorporating cytoskeletal, extracellular-matrix, confinement, and phenotype mechanisms, while reporting no validated consensus on the essential subcellular processes or the minimal adequate detail. These are contextual mechanisms mentioned in the curation, not a prescribed list of essential subcellular processes. The record labels the problem open; this label and the assessment are supplied dated curation, not independently verified mathematical conclusions.
\subsection{Short English statement}
Determine which subcellular processes must be represented to explain cancer-cell invasion, and how much mechanistic detail a realistic, sufficiently simple model needs. The source leaves the biological setting, meaning of necessity, validation criteria, and measure of simplicity unspecified, so a unique mathematical formulation requires further choices. This is a statement-only draft requiring human mathematical review.
\subsection{Sources}
\begin{itemize}
\item Supplied local UnsolvedMath registry record, permanent ID \(30002570\), title ``Minimal Mechanistic Models of Cancer-Cell Invasion'', source problem code \texttt{catalogue entry}, in the collection ``Oberwolfach Reports Open Problems''. No individual proposer is identified. The record gives the proposed year as 2014.
\item Supplied report citation: \emph{Mini-Workshop: Mathematical Models for Cancer Cell Migration (2015)}, Oberwolfach Reports, DOI 10.4171/owr/2014/19, pp.\ 1075--1109. \url{https://doi.org/10.4171/owr/2014/19}. The background describes Oberwolfach Reports as documenting mathematical workshops and their open-problem discussions. The citation's 2015 designation and the record's 2014 proposed year are retained without resolving their relationship.
\item Supplied literature source, ``Multiscale invasion model'': \url{https://pmc.ncbi.nlm.nih.gov/articles/PMC10293590/}.
\item Supplied literature source, ``Molecular-pathway invasion model'': \url{https://doi.org/10.1371/journal.pcbi.1004571}.
\item Literature assessment checked on \texttt{2026-08-21}, according to the supplied record: status ``open''; multiscale models incorporate cytoskeletal, extracellular-matrix, confinement, and phenotype mechanisms, but the assessment reports no validated consensus identifying essential subcellular processes or a minimal adequate level of detail. This is supplied dated literature triage derived from the OWR curation, not independently verified here. Time-sensitive status and cited claims require independent verification.
\item The registry marks the statement status ``exact'' and supplies a verification note describing it as clear, self-contained, publication-ready, and retained byte-for-byte. These are source-record assertions, not an independent verification of the report wording or of mathematical completeness.
\end{itemize}
\end{document}
